\section{Related work}
\label{sec:related}

\paragraph{DEM inversion.}
Classical approaches include Markov chain Monte Carlo~\cite{kashyap1998mcmc}, regularized inversion~\cite{hannah2012regularized}, fast approximate schemes~\cite{plowman2013fast,plowman2020simple}, sparse inversion by Basis Pursuit~\cite{cheung2015dem,chen1998basis} and Elastic Net~\cite{athiray2024enet,zou2005elastic}. We take the last two as given and ask how to amortize them without running them.

\paragraph{Neural emulators for solar data.}
Szenicer \etal~\cite{szenicer2019virtual} train deep networks on AIA images to predict irradiance, and SuperSynthIA~\cite{wang2024supersynthia} produces full-disk magnetograms from per-pixel distribution outputs. For DEMs, DeepEM~\cite{wright2019deepem} regresses BP's output. DEMNet~\cite{perezdiaz2026demnet} does the same at scale and then finetunes on a forward-model term, and Balodhi and Morton~\cite{bayesdem2026} mix a forward-model term with a loss on regularized-inversion DEMs and estimate uncertainty with dropout. All three need solver DEMs. DEMNet notes that a forward-model term is hard to optimize from scratch because negative solutions also minimize it. We avoid this by predicting the solver's non-negative weights and training on its whole objective, which fixes the tie-break as well as the fit. To our knowledge no earlier DEM network trains without solver labels.

\paragraph{Amortized and constrained optimization.}
Amortized optimization replaces repeated solves with a learned map~\cite{amos2023tutorial}. LISTA~\cite{gregor2010lista} did this for sparse coding by regressing solver output. Closer to us are methods that train on the objective itself. Ablin \etal~\cite{ablin2019unfolded} train unfolded ISTA on the Lasso cost, and other work learns constrained programs without solver labels through completion and correction~\cite{donti2021dc3}, Lagrangian duality~\cite{fioretto2020lagrangian}, primal-dual training~\cite{park2023primal}, feasibility layers~\cite{chen2024e2elr} and exact penalties~\cite{benslimane2025penalty}; Kotary \etal~\cite{kotary2021survey} give a survey. Those methods mostly aim at feasibility on structured problem families. Our instances are tiny but numerous, observation noise makes a small violation harmless, and the question is how large that violation is at a given weight. The relevant theory is the classical penalty method~\cite{courant1943variational,fiacco1968nonlinear,nocedal2006numerical}, usually read as a convergence result. Mangasarian~\cite{mangasarian2004newton} showed that a finite penalty on an LP already recovers an exact dual solution, and we use that fact as a calibration. Physics-informed networks~\cite{raissi2019pinn} use similar weighted penalties and show the same training stiffness as the weight grows~\cite{wang2021pathologies,krishnapriyan2021failure}.

\paragraph{Forward models in vision and astronomy.}
Deep inverse-problem methods~\cite{ongie2020inverse} often fit one scene through a differentiable forward model, from deep image priors~\cite{ulyanov2018dip} to neural fields for black-hole emission~\cite{levis2022bhnerf,feng2026pidef} and the cosmic web~\cite{zhao2025cosmic}. We share the forward-model loss but amortize over pixels and images, so no optimization runs at test time. Our analysis of blow-ups connects to work on how networks behave far from their training data~\cite{hein2019relu,xu2021extrapolate}.
